% !TeX program = lualatex
% !TeX encoding = utf8
% !TeX spellcheck = uk_UA
% !TeX root =../ClassicalEletrodynamics.tex

%=========================================================
\chapter{Перетворення Лоренца як наслідок постулатів Айнштайна}\label{\currfilebase}
%=========================================================

У цьому розділі викладено основні поняття спеціальної теорії відносності (СТВ) та їх найбільш важливі наслідки, що охоплюють діапазон від субатомних до астрономічних  масштабів. З базових принципів СТВ виведено перетворення Лоренца згідно до оригінальної роботи Айнштайна 1905 р., показано інваріантність квадрату інтервалу, розглянуто перетворення довжини та об'єму.

%% --------------------------------------------------------
\section{Простір-час, системи відліку}
%% --------------------------------------------------------

Спеціальна теорія відносності (СТВ) переглянула фундаментальні поняття класичної фізики. Щоб з'ясувати, у чому полягає цей перегляд, необхідно відштовхуватися від найбільш елементарних уявлень про вимірювання часу й відстаней. Виявляється, що не усі припущення ньютонівської теорії, зокрема, стосовно існування абсолютного часу та поняття одночасності, наповнені фізичним змістом.

Одним з вихідних у теорії відносності є поняття події, яку характеризують однією часовою та трьома просторовими координатами у певній системі відліку. Це елементарне, базове поняття, яке розкривається через взаємовідношення з іншими базовими поняттями, так само, як поняття точки в геометрії. Уявлення про простір-час як множину подій, яка має властивості многовиду, дозволяє розглядати в ньому кількісні співвідношення й використовувати потужний апарат математичного аналізу і диференціальної геометрії.

Будемо говорити, що визначена система відліку (СВ), якщо відома фізична процедура, що дозволяє однозначно зіставити будь-якій події четвірку координат --- часову й три просторові координати. Практично СВ реалізують за допомогою певних технічних засобів та алгоритмів, що дозволяють визначати координати за даними спостережень, як, наприклад, в астрометричних СВ --- геоцентричній та геліоцентричній тощо. Поняття СВ співзвучне поняттю системи координат, але відрізняється від нього, оскільки передбачає певні, інколи ідеалізовані, фізичні дії, вимірювання тощо. У класичній (тобто, неквантовій) теорії поля припускають, що ці вимірювання можуть бути виконані як завгодно точно, хоча на практиці вони підлягають певним обмеженням і виконуються з певною точністю.

Далі ми будемо вважати, що простір і час є однорідними, а простір є ізотропним. Однорідність простору і часу означає, що результати будь-яких фізичних експериментів не залежать ні від розташування лабораторії, ані від часу проведення. Ізотропія означає таку незалежність від просторової орієнтації лабораторії. Однорідність простору і часу та ізотропія простору --- це певна ідеалізація, яка порушується, зокрема, за наявності гравітаційних полів, що вивчає загальна теорія відносності (ЗТВ). Задача фізика --- визначити, наскільки ця ідеалізація можлива в умовах конкретного експерименту.

Наступна ідеалізація --- інерціальні системи відліку, тобто такі, в яких виконується закон інерції: декартові координати матеріальної точки, на яку не діють зовнішні сили (або сума сил дорівнює нулю), є лінійними функціями часу. Ми припускаємо, що інерціальні системи існують, а просторові співвідношення в інерціальних СВ відповідають аксіомам евклідової геометрії. Евклідова аксіоматика (знов-таки, за відсутності гравітаційних ефектів, які вивчає ЗТВ) підтверджується високоточними експериментами. В кожній інерціальній СВ визначені інтервали часу та відстані між будь-якими двома подіями, і відповідні вимірювальні процедури не призводять до суперечностей. Евклідовість простору дозволяє, зокрема, застосовувати прямокутні декартові системи координат в інерціальних СВ, де квадрат довжини відрізка прямої дорівнює, за теоремою Піфагора, сумі квадратів проекцій відрізка на координатні осі.

Інерціальну систему відліку в СТВ можна уявити як сукупність спостерігачів з годинниками, що щільно заповнюють простір і є нерухомими один відносно одного, причому годинники синхронізовані. Нерухомість спостерігачів один відносно одного можна перевірити, вимірюючи час від часу затримку сигналів, що генерують передавачі біля кожного спостерігача. Це також дає змогу визначити відстань між спостерігачами, знаючи швидкість світла у вакуумі $c$, яка є однаковою в усіх інерціальних системах --- це один з базових постулатів СТВ, який перевірено численними експериментами. Але це не означає, що числове значення відстані між спостерігачами буде одне й те саме в різних СВ, що рухаються одна відносно іншої.

Спеціальна теорія відносності розглядає передусім усі СВ, які вважаємо інерціальними. При розгляді перетворень між інерціальними системами, які розглянуто далі, також зазвичай вважаємо, що просторові координати подій в інерціальних СВ визначені в прямокутній декартовій системі.

Однакові нерухомі відносно деякої СВ годинники вважаються синхронізованими, якщо в них визначений спільний початок відліку. Синхронізацію в заданій СВ можна провести, наприклад, за допомогою електромагнітних сигналів. Розглянемо одну з можливих процедур синхронізації (за Айнштайном). Нехай події $A$ та $B$ супроводжуються світловими спалахами, а спостерігач, нерухомий відносно СВ, що знаходиться на прямій посередині між просторовими точками подій, реєструє ці спалахи. Тоді події $A$ та $B$ є одночасними в даній СВ, якщо цей спостерігач реєструє сигнали від спалахів одночасно.

Розглянемо інший, еквівалентний, спосіб синхронізації за допомогою світлових сигналів. Нехай спостерігачі А та В зі своїми годинниками нерухомі відносно деякої СВ. Зафіксуємо початок відліку годинника А та зкорегуємо відліки годинника В, щоб можна було вважати його синхронізованим з А. Щоб зіставити покази годинників А і В, спостерігач А у точці $x_A$ посилає світловий сигнал в момент $t_A$ (за своїм годинником). Цей сигнал відбивається від дзеркала біля спостерігача В (подія $B$). Спостерігач А реєструє відбитий сигнал у своєму околі в момент $t_A'$. Позначимо $t_B = \frac{t_A + t_A'}{2}$. Будемо вважати, що, за означенням, подія $B$ є одночасною з подією $A$. Досі тут ніде не фігурували покази годинника В. Але тепер можна встановити початок відліку годинника В так, щоб події $B$ за його показами також відповідав час $t_B$. Описана процедура дає змогу встановити однаковий початок відліку годинників усіх спостерігачів, нерухомих один відносно одного. Підкреслимо, що процедура синхронізації стосується тільки однієї СВ.

Наприкінці відзначимо важливі відмінності СТВ від ньютонівської теорії:
\begin{itemize}
\item[(i)] в СТВ події, що є одночасними в одній інерціальній СВ, можуть не бути одночасними в іншій інерціальній СВ;
\item[(ii)] якщо два спостерігача рухаються один відносно одного, їх годинники неможливо синхронізувати.
\end{itemize}
Це суттєво відрізняє СТВ від ньютонівського підходу з абсолютним часом, загальним для усіх систем відліку. Але саме СТВ узгоджується з фізичним досвідом.

%% --------------------------------------------------------
\section{Перетворення Лоренца для одновимірних рухів систем відліку}
%% --------------------------------------------------------

Перейдемо до перетворень, що пов'язують координати події $(x, t)$ та $(x', t')$ в двох інерціальних системах $S$ та $S'$. Як зазначено вище, вважаємо, що просторові координати $x, y, z$ та $x', y', z'$ в $S$ та $S'$ визначені в прямокутній декартовій системі. Нехай осі абсцис в $S$ та $S'$ збігаються. Для початку обмежимось одновимірним випадком, коли події відбуваються на осях $x$, тобто перетворюються лише $x$ і $t$ ($y=z=0$, $y'=z'=0$).

\begin{equation}
x = x(x', t'), \quad t = t(x', t')
\label{eq:1.1}
\end{equation}

Для малих змін координат $dx$ та $dt$
\[
dx = \frac{\partial x}{\partial x'} dx' + \frac{\partial x}{\partial t'} dt', \quad
dt = \frac{\partial t}{\partial x'} dx' + \frac{\partial t}{\partial t'} dt'.
\]

Величини $\frac{\partial x}{\partial x'}$ тощо можна вимірювати відповідно в $S$ та $S'$. Це дає принципову можливість експериментально визначити частинні похідні від $x$ та $t$. Наприклад, розглядаючи дві події з координатами $(x_1, t)$ та $(x_2, t)$ в $S$, одночасні в цій системі, та вимірюючи $x_1'$ та $x_2'$ в $S'$, з рівнянь для диференціалів визначимо $\frac{\partial x}{\partial x'}$ та $\frac{\partial x}{\partial t'}$ і так далі. Результати такого визначення можуть залежати від відносної швидкості систем відліку, але не повинні залежати від положення та часу вимірювань згідно з гіпотезою про однорідність простору-часу. Тому ці похідні є константами і перетворення (\ref{eq:1.1}) є лінійним. Запишемо його у вигляді

\begin{equation}
x = a_{11} x' + a_{12} t'
\label{eq:1.2}
\end{equation}

\begin{equation}
t = a_{21} x' + a_{22} t'
\label{eq:1.3}
\end{equation}
де $a_{ij}$ --- константи.

Очевидно, цілком аналогічні міркування доводять лінійність і в більш загальному, ніж одновимірний, випадку. Лінійність перетворень, що пов’язують координати подій в інерціальних СВ, є наслідком однорідності простору-часу.

Проаналізуємо перетворення \eqref{eq:1.2}, \eqref{eq:1.3}, постулюючи принцип відносності:



% Використання:
\begin{postulate}{}
Усі інерціальні системи відліку рівноправні.
\end{postulate}

Це означає, що всі фізичні процеси виглядають однаково для різних інерціальних спостерігачів, що рівномірно рухаються зі сталими швидкостями і знаходяться в різних просторово-часових областях, тобто усі процеси в різних інерціальних системах відліку можна описати однаковими рівняннями незалежно від відносної швидкості цих систем. Неможливо встановити ``абсолютну'' СВ, відносно якої загальні фізичні закони мали б якусь специфіку, якої немає в інших СВ. З точки зору принципу відносності, який узгоджується з усіма спостереженнями, сама постановка задачі про пошук ``абсолютної'' СВ беззмістовна, як і не мають фізичного сенсу поняття про абсолютний спокій чи рух.

Сформульований принцип дуже загальний і для просування далі необхідна додаткова інформація, яка, зокрема, може міститися в рівняннях динаміки полів чи частинок. Сам по собі принцип відносності можна узгодити навіть з ньютонівськими уявленнями про єдиний час для усіх інерціальних систем, при цьому, якщо рівняння ньютонівської динаміки були б правильними, тоді б перетворення, що пов'язують різні системи відліку, мали вигляд перетворень Галілея, а не перетворень Лоренца, що будуть виведені далі. Додаткову інформацію, що конкретизує зв'язок інтервалів часу та відстаней найбільш простим чином, дає другий постулат --- умова інваріантності швидкості світла в усіх інерціальних системах:

\begin{postulate}{}
Швидкість світла у вакуумі є універсальною сталою, вона однакова в усіх інерціальних системах відліку і не залежить від швидкості джерела.
\end{postulate}

Як побачимо далі, спираючись на цей постулат разом із принципом відносності та гіпотезою про однорідність простору-часу та ізотропію простору, можна отримати зв'язок координат в різних інерціальних системах, не звертаючись до рівнянь Максвелла чи якихось інших рівнянь теорії поля. Спочатку перетворення, що описують цей зв'язок, з'явилися в роботах Лоренца та Пуанкаре саме на основі аналізу рівнянь Максвелла. Однак саме Айнштайн запропонував розглядати принцип відносності разом з вимогою інваріантності швидкості світла і надав цим твердженням фундаментальний зміст. Постулати відносності та інваріантності швидкості світла, якщо їх розглядати разом, відомі як спеціальний принцип відносності Айнштайна.

Твердження про інваріантність швидкості світла заперечує існування виділеної системи відліку, пов'язаної з ефіром --- особливим середовищем, збудження якого обумовлюють поширення електромагнітних хвиль. За уявленнями, що панували на початку 20 століття до створення СТВ, система спокою ефіру займає виділене положення, час, визначений в цій ``абсолютній'' СВ, є однаковим в усіх інерціальних СВ, а події, що є одночасними в одній інерціальній системі, є одночасними в усіх таких системах. Але гіпотеза про існування ефіру була спростована багатьма експериментами. Ньютонівська механіка не узгоджується зі спостереженнями, хоча для швидкостей, малих у порівнянні зі швидкістю світла, ця механіка працює досить точно, а відповідні релятивістські поправки (тобто поправки СТВ) до ньютонівських розрахунків спостережуваних величин мають порядок $v^2/c^2$. Щодо рівнянь Максвелла, які самі по собі не вимагають наявності ефіру, то вони чудово узгоджуються зі спеціальним принципом відносності.

Використаємо спеціальний принцип відносності для встановлення явного вигляду координатних перетворень (\ref{eq:1.2}), (\ref{eq:1.3}), що пов'язують інерціальні системи. Обмежимось однорідними перетвореннями, коли подія з координатами $(x', t')$ в $S'$ має також координати $(x, t)$ в $S$, тобто $x=0$, $t=0$ при $x'=0$, $t'=0$. Цього можна завжди домогтися перевизначенням початку відліку часу та просторових координат.

Нехай, наприклад, початок координат системи $S'$ віддаляється від початку координат $S$ зі швидкістю $v$, тобто $S'$ рухається в додатному напрямку осі $x$.

Якщо в $S'$ траєкторія деякої точки, що покоїться в цій системі, визначається співвідношенням $x' = 0$ при усіх $t'$, то в $S$ маємо $x = vt$. Звідси з (\ref{eq:1.2}) маємо $a_{12} = v a_{22}$, і замість (\ref{eq:1.2}) можна записати

\begin{equation}
x = \gamma (x' + v t')
\label{eq:1.4}
\end{equation}

Аналогічно, сукупність подій з координатою $x = 0$ в $S$ (траєкторія частинки, що покоїться відносно $S$) відповідає траєкторії $x' = -v t'$ в $S'$; звідси визначаємо $a_{21} = -v a_{11}$ і замість (\ref{eq:1.3}) маємо

\begin{equation}
t = \gamma' \left(t' - \frac{v}{c^2} x'\right)
\label{eq:1.5}
\end{equation}

Врахуємо, що швидкість світла в $S$ та $S'$ однакова. Це означає, що траєкторія світлового сигналу, задана рівнянням $x = c t$ в $S$, матиме такий самий вигляд $x' = c t'$ в $S'$.

Підставимо це в (\ref{eq:1.4}), (\ref{eq:1.5}), звідки отримаємо $\gamma = \gamma'$.

Таким чином, для будь-якої події $(x', t')$ її координати в $S$ визначаються з формул

\begin{equation}
x = \gamma (x' + v t'), \quad t = \gamma \left(t' + \frac{v}{c^2} x'\right).
\label{eq:1.6}
\end{equation}

Обернений зв'язок отримуємо, розв'язуючи ці рівняння відносно $x$, $t$:

\begin{equation}
x' = \gamma (x - v t), \quad t' = \gamma \left(t - \frac{v}{c^2} x\right),
\label{eq:1.7}
\end{equation}
де

\begin{equation}
\gamma = \frac{1}{\sqrt{1 - \frac{v^2}{c^2}}}.
\label{eq:1.8}
\end{equation}

Зауважимо, що $\gamma > 0$, оскільки перетворення (\ref{eq:1.6}) повинні мати загальний характер, а перехід від координат в $S$ до $S'$ відповідає оберненому знаку швидкості $v$.

Поки що невідомо, чи $\gamma$ збігається з $\gamma'$. Щоб це з’ясувати, згадаємо припущення про ізотропію простору. Поміняємо напрямок просторових осей ($v \to -v$ в $S$ та $S'$), тобто вводимо системи $S''$ з координатами $(x'', t'')$ та систему $S'''$, де $v \to -v$. Взаємне розташування $S''$ відносно $S$ таке ж як і $S'$ відносно $S$ (початок координат $S''$ віддаляється від початку координат $S$ так само, як і $S'$ відносно $S$), тому вигляд координатних перетворень $S \to S''$ має бути такий самий, як (\ref{eq:1.6}).

Підставляючи $x'' = x'$, $t'' = t'$ та змінюючи знак $v$, отримуємо перехід від $S'$ до $S$.

Порівнюючи це з (\ref{eq:1.7}), маємо $\gamma = \gamma'$, звідки, враховуючи (\ref{eq:1.8}),

\begin{equation}
\gamma = \frac{1}{\sqrt{1 - \frac{v^2}{c^2}}}.
\label{eq:1.9}
\end{equation}

Формально тут є ще одна можливість, коли $\gamma = -1/\sqrt{1 - v^2/c^2}$. Цей варіант перетворення координат є добутком перетворення з множником (\ref{eq:1.9}) та інверсії часової та просторової координат. Нас цікавить, однак, неперервна залежність $\gamma(v)$, причому маємо отримати тотожнє перетворення при $v=0$. Тому варіант з від'ємним $\gamma$ відкидаємо.

Дамо більш наочний аналіз умови ізотропії на прикладі процесу випромінювання, що виникає внаслідок квантових переходів в атомі водню, з точки зору спостерігачів, пов'язаних з $S$ та $S'$. Нехай відносна швидкість систем $v$. Отримаємо зв'язок між $\Delta t$ і $\Delta t'$ виходячи з (\ref{eq:1.6}),(\ref{eq:1.7}). Відраховуючи певну кількість періодів коливань квантової системи, можна створити атомний годинник (еталон часу), який генерує спалахи світла з власним періодом $\tau_0$. Нехай цей годинник та спостерігач поблизу нього, що реєструє спалахи, покояться в системі $S'$ у точці $x'=0$. Розглянемо дві події реєстрації спалахів з координатами $(0, t_1')$ в $S'$ та $(0, t_2')$. Для спостерігача в іншій системі $S$ (з таким самим годинником) інтервал часу між цими подіями визначається згідно (\ref{eq:1.7}).

Таким чином, відношення $\Delta t / \Delta t'$ може бути визначено експериментально.

Розглянемо тепер аналогічне вимірювання в $S$ та $S''$, але нехай тепер годинник покоїться в $S$ і інтервал $\Delta t$ в цій системі визначають тією самою кількістю коливань, що й раніше. Цей експеримент відрізняється від попереднього лише напрямком швидкості між системами, яка, за вимогою ізотропії, несуттєва. Враховуючи (\ref{eq:1.6}), аналогічно матимемо $\Delta t'' / \Delta t = \gamma$, тобто $\Delta t' / \Delta t = 1/\gamma$.

Запишемо остаточний результат

\begin{equation}
x = \gamma (x' + v t'), \quad t = \gamma \left(t' + \frac{v x'}{c^2}\right),
\label{eq:1.10}
\end{equation}

а також обернене перетворення

\begin{equation}
x' = \gamma (x - v t), \quad t' = \gamma \left(t - \frac{v x}{c^2}\right).
\label{eq:1.11}
\end{equation}

Прямою підстановкою легко перевірити, що внаслідок (\ref{eq:1.10}) або (\ref{eq:1.11})

\[
x^2 - c^2 t^2 = x'^2 - c^2 t'^2.
\]

Досі в п.1.2 розглядалися лише події, що відбуваються на осях абсцис систем $S$ та $S'$. Але усі міркування легко поширити на довільні події; однак відносну орієнтацію систем $S$ та $S'$, просторові осі яких паралельні, поки що залишимо як раніше. Завдяки інваріантності швидкості світла та ізотропії простору, фронт світлової хвилі повинен мати однакову сферичну форму в обох системах $S$ та $S'$:

\[
x^2 + y^2 + z^2 - c^2 t^2 = x'^2 + y'^2 + z'^2 - c^2 t'^2,
\]

звідки $y' = y$, $z' = z$ для довільних $x, t$. Звідси випливає, що координати, осі яких перпендикулярні напрямку відносного руху, залишаються незмінними:

\begin{equation}
y' = y, \quad z' = z,
\label{eq:1.11a}
\end{equation}

де ми, аналогічно попередньому, відкидаємо варіант з інверсією просторових осей. Перетворення (\ref{eq:1.10}), (\ref{eq:1.11}) носять ім’я Лоренца, який отримав їх, вивчаючи інваріантність рівнянь Максвелла. Далі називатимемо перетворення (\ref{eq:1.10}), (\ref{eq:1.11}) одновимірними перетвореннями Лоренца, щоб відрізняти їх від загальних перетворень, що пов'язують інерціальні системи відліку з довільною орієнтацією просторових осей та довільним напрямком відносного руху.

\textbf{Відносність одночасності.} Покажемо, що одночасні події в $S$ з координатами $(x_1, t)$, $(x_2, t)$, які відбуваються у різних просторових точках $x_1 \neq x_2$, не є одночасними в $S'$. Підставляючи координати подій в (\ref{eq:1.10}), маємо

\[
t_1' = \gamma \left(t - \frac{v x_1}{c^2}\right), \quad t_2' = \gamma \left(t - \frac{v x_2}{c^2}\right).
\]

Таким чином, події з різними значеннями координати $x$, одночасні в одній інерціальній системі $S$, не є одночасними в іншій системі $S'$, що рухається відносно першої вздовж $x$ з ненульовою швидкістю $v$. Однак одночасність зберігається для подій з однаковими $x$, але різними координатами $y$ або $z$. Дійсно, в перетворення часової змінної $t'$ чи $t$ не входять $y$ і $z$. Тому, наприклад, одночасні події на осі $y$ в $S$ є одночасними і в $S'$.

\textbf{Додавання швидкостей.} Перетворення Лоренца пов'язують координати будь-яких подій в інерціальних СВ $S$ та $S'$, де $S'$ рухається відносно $S$ із сталою швидкістю $v$. Якщо розглядається рух частинки вздовж осі $x$ системи $S'$ з деякою швидкістю $u' = dx'/dt'$, то в системі $S$, яка рухається відносно $S'$ зі швидкістю $v$, маємо

\[
u = \frac{dx}{dt} = \frac{\gamma (dx' + v \, dt')}{\gamma (dt' + \frac{v}{c^2} dx')} = \frac{u' + v}{1 + \frac{u' v}{c^2}}.
\]

Аналогічно

\[
u' = \frac{u - v}{1 - \frac{u v}{c^2}}.
\]

Це --- формула релятивістського додавання швидкостей у випадку руху вздовж однієї осі.

Нехай тепер деяке тіло А рухається зі швидкістю $u$ відносно спостерігача, а тіло В --- зі швидкістю $u'$ відносно А (в одному напрямку). Рухи вважаємо інерціальними. За формулою релятивістського додавання швидкостей тіло В рухається відносно спостерігача зі швидкістю

\[
w = \frac{u + u'}{1 + \frac{u u'}{c^2}}.
\]

Розглянемо добуток

\[
\frac{w}{c} = \frac{\frac{u}{c} + \frac{u'}{c}}{1 + \frac{u u'}{c^2}}.
\]

Звідси видно, що в разі одновимірних рухів величина $\frac{u}{c}$ є адитивною, подібно для звичайної швидкості в ньютонівській теорії.

%% --------------------------------------------------------
\section{Власний час}
%% --------------------------------------------------------

Нехай точкове тіло рухається зі сталою швидкістю $u$ вздовж осі $x$. Розглянемо СВ $S'$, яка рухається разом із тілом, тобто в якій тіло покоїться. Таку систему називають власною для даного тіла. Розглянемо дві послідовні події на траєкторії частинки $1$ і $2$, де $\Delta \tau$ --- інтервал власного часу, тобто часу, що вимірюється годинниками у власній СВ. Величина $\Delta \tau$, очевидно, не залежить, відносно якої СВ ми розглядаємо рух частинки.

Нехай у системі $S$ координати цих подій є $(x_1, t_1)$, $(x_2, t_2)$, а тіло рухається зі швидкістю $u$. Записуючи зв’язок між координатами подій в $S$ і $S'$ згідно до (\ref{eq:1.11}), отримуємо

\begin{equation}
\Delta \tau = \Delta t \sqrt{1 - \frac{u^2}{c^2}}.
\label{eq:1.12}
\end{equation}

Ця формула дає змогу обчислювати власний час по траєкторії частинки, що задана в будь-якій інерціальній СВ.

Виникає питання: чи застосовна формула (\ref{eq:1.12}) у випадку неінерціальних рухів? Цей результат не випливає з СТВ. Адже якби формула для власного часу містила поправку, пропорційну прискоренню, яка зникає у випадку інерційних рухів, це не суперечило б наведеним вище міркуванням, що привели до (\ref{eq:1.12}). Відповідь дають численні експерименти, які показують, що власний час (\ref{eq:1.12}) не залежить від прискорення. Тому для будь-якої ділянки траєкторії частинки $x(t)$, її власний час обчислюється за формулою

\begin{equation}
\Delta \tau = \int \sqrt{1 - \frac{u^2(t)}{c^2}} \, dt, \quad u(t) = \frac{dx}{dt},
\label{eq:1.13}
\end{equation}

яка застосовна також у випадку неінерціальних рухів. Підкреслимо, що в (\ref{eq:1.13}) рух частинки може бути прискореним, але траєкторія частинки має описуватися в інерціальній СВ. В неінерціальних СВ формулу (\ref{eq:1.13}) слід модифікувати.

Переписуючи (\ref{eq:1.13}) для диференціалів, маємо

\[
d\tau = \frac{1}{c} \sqrt{c^2 dt^2 - dx^2 - dy^2 - dz^2}.
\]

Цю величину називають квадратом інтервалу; тут він записаний для подій, що мають місце на траєкторії деякої частинки, швидкість якої менша за швидкість світла.


%% --------------------------------------------------------
\section{Квадрат інтервалу}
%% --------------------------------------------------------

Узагальнимо твердження про інваріантність квадрата інтервалу на довільні відносні рухи інерціальних систем. Нехай $(x, y, z, t)$ --- координати будь-яких подій в СВ $S$, а $(x', y', z', t')$ --- відповідні координати в $S'$, яка рухається в напрямку осі $x$ відносно $S$. Тоді за формулами (\ref{eq:1.10}),

\begin{equation}
ds^2 = ds'^2.
\label{eq:1.14}
\end{equation}

У більш загальному випадку, коли $S'$ рухається відносно $S$ зі швидкістю $\mathbf{v}$ у довільному напрямку, розглянемо події з координатами $(x,y,z,t)$ в $S$ та відповідними їм $(x',y',z',t')$ в $S'$. Порівняємо в обох системах квадрат інтервалу, який, за визначенням, дорівнює

\begin{equation}
ds^2 = c^2 dt^2 - dx^2 - dy^2 - dz^2, \quad
ds'^2 = c^2 dt'^2 - dx'^2 - dy'^2 - dz'^2.
\label{eq:1.15}
\end{equation}

Очевидно, $ds^2$, а значить і $ds'^2$, не змінюється при просторових обертаннях декартової системи координат. Таким чином, величина квадрату інтервалу є незмінною для довільних відносних рухів інерціальних систем відліку: $ds^2 = ds'^2$.

Кажуть, що дві події пов'язані часоподібним інтервалом, якщо $ds^2 > 0$, та просторовоподібним інтервалом, якщо $ds^2 < 0$. Легко бачити, що нескінченно близькі події вздовж 4-траєкторії частинки (світової лінії) пов'язані часоподібним інтервалом, оскільки

\[
ds^2 = c^2 d\tau^2 > 0
\]
%
внаслідок умови $u < c$. Відповідно, таку траєкторію називають часоподібною. Якщо частинка рухається зі швидкістю $u = c$, маємо $ds^2 = 0$ і траєкторію називають світлоподібною чи ізотропною.

%% --------------------------------------------------------
\section{Перетворення довжини та об'єму}
%% --------------------------------------------------------

Повернемося знов до систем, що пов’язані перетвореннями (\ref{eq:1.10}), (\ref{eq:1.11}). Нехай $S'$ є власною системою для деякого стержня довжиною $L_0$. Порівняємо довжину стержня в різних системах, вважаючи, що він орієнтований вздовж осі абсцис.

Стержень рухається відносно лабораторної системи $S$, тому події, що відповідають вимірюванням координат кінців стержня в цій системі, мають відбуватися одночасно. Позначимо відповідні координати $x_1, x_2, t$, $L = x_2 - x_1$ --- довжина стержня в $S$. У власній системі кінці стержня мають весь час постійні координати. Просторові координати подій в $S'$

\[
x_1' = \gamma (x_1 - v t), \quad x_2' = \gamma (x_2 - v t),
\]

звідси

\begin{equation}
L = \frac{L_0}{\gamma}.
\label{eq:1.16}
\end{equation}

Обчислення нагадують виведення формули для інтервалів часу

\begin{equation}
\Delta \tau = \frac{\Delta t}{\gamma},
\label{eq:1.17}
\end{equation}
%
однак у випадку (\ref{eq:1.16}) розглядаються події з однаковими часовими координатами в лабораторній системі, відносно якої стержень рухається, а при виведенні (\ref{eq:1.17}) розглядаються події з однаковими просторовими координатами у власній системі годинника.

Нехай $S'$ з координатами $x',y',z',t'$ є власною системою для паралелепіпеда

\[
\Delta x' \Delta y' \Delta z' = \Delta V_0;
\]

відповідний елемент об’єму $\Delta V_0$. Враховуючи, що $\gamma = 1/\sqrt{1 - v^2/c^2}$, маємо в лабораторній системі, відносно якої об'єм рухається зі швидкістю $v$,

\begin{equation}
\Delta V = \frac{\Delta V_0}{\gamma}.
\label{eq:1.18}
\end{equation}

В разі зарядженого середовища, отримаємо зв'язок густин заряду в різних інерціальних системах. Нехай $\rho_0$ --- густина заряду у власній системі, відносно якої заряди покояться, $\Delta V_0$ --- об’єм в цій системі з сумарним зарядом $Q$, $\Delta V$ --- об’єм цього ж заряду у системі, відносно якої заряд рухається з швидкістю $v$. Нагадаємо, що заряд має одну й ту ж саму величину в різних системах. Це дає

\begin{equation}
\rho = \gamma \rho_0.
\label{eq:1.19}
\end{equation}

Звідси випливає інваріантність величини $Q$.

%=========================================================
% ВПРАВИ
%=========================================================

%=========================================================
\begin{problem}
Показати прямим обчисленням, що добуток одновимірних перетворень Лоренца, що відповідають швидкостям $v_1, v_2$, є також перетворення Лоренца з новою швидкістю $v = \frac{v_1+v_2}{1+v_1 v_2/c^2}$.
\end{problem}
%=========================================================

%=========================================================
\begin{problem}
Відносно інерціальної системи $K$ з координатами $(t, \mathbf{r})$ рухається із швидкістю $\mathbf{v}$ інша інерціальна система $K'$ з координатами $(t', \mathbf{r}')$. Записати відповідне перетворення $\mathbf{r}'(\mathbf{r}, t), t'(\mathbf{r}, t)$. Напрямок тривимірного вектора $\mathbf{v}$ довільний.
Ввказівка: Тривимірний вектор $\mathbf{r}$ розбити на поздовжню та поперечну відносно $\mathbf{v}$ частини: $\mathbf{r} = \mathbf{r}_{||} + \mathbf{r}_{\perp}$ (поздовжня перетворюється за одновимірним перетворенням Лоренца, поперечна залишається незмінною).
\end{problem}
%=========================================================

%=========================================================
\begin{problem}
Показати прямим обчисленням, що перетворення координат з попередньої вправи зберігає квадрат інтервалу $s^2 = c^2t^2 - \mathbf{r}^2$.
\end{problem}
%=========================================================

%=========================================================
\begin{problem}
Знайти закон перетворення прискорення для одновимірних рухів при переході в систему $K'$, що рухається зі швидкістю $v$ відносно системи $K$.
Відповідь: $a' = \frac{a(1-v^2/c^2)^{3/2}}{(1-uv/c^2)^3}$.
\end{problem}
%=========================================================

%=========================================================
\begin{problem}
Швидкість ракети, що рухається вздовж прямої, завжди менша за $c$. Чи можлива з точки зору СТВ ситуація, коли радіосигнал, що поширюється за ракетою у тому ж напрямку, не наздоганяє її? Розгляньте приклад, коли космонавт у ракеті має завжди сталу вагу, як на Землі; за початкові умови візьміть $x(0)=c^2/g, u(0)=0$.
\end{problem}
%=========================================================

%=========================================================
\begin{problem}\label{prb:abberation}
Спостерігач вимірює кут $\theta$ між зірками $A$ та $B$. Знайти кут $\theta'$, який вимірює космонавт, що рухається зі сталою швидкістю $v$ відносно спостерігача в напрямку зірки $B$.

\emph{Відповідь}: $\cos\theta' = \frac{\cos\theta - v/c}{1 - (v/c)\cos\theta}$.
\end{problem}
%=========================================================

%=========================================================
\begin{problem}
Розподіл далеких зірок на небесній сфері в системі спостерігача є ізотропним. Знайти розподіл зірок $\frac{dN}{d\Omega'}$, який бачить космонавт, що рухається зі швидкістю $v$ відносно спостерігача.

\textbf{Відповідь:} $\frac{dN}{d\Omega'} = \frac{N}{4\pi} \frac{1-v^2/c^2}{(1 + \frac{v}{c}\cos\theta')^2}$, де $N$ — повне число зірок, $d\Omega' = \sin\theta' d\theta' d\phi'$.
\end{problem}
%=========================================================

%=========================================================
\begin{problem}
Розподіл далеких зірок в системі спостерігача є ізотропним. Космонавт рухається зі сталою швидкістю $v$ відносно спостерігача. Знайти відношення числа зірок, що спостерігає космонавт у передній напівсфері, до числа зірок у задній напівсфері.

\emph{Вказівка:} скористатися безпосередньо результатом вправи \ref{prb:abberation}.

\emph{Відповідь:} $\frac{1+v/c}{1-v/c}$.

\end{problem}
%=========================================================

%=========================================================
\begin{problem}
Розпад $\pi^0 \to 2\gamma$ відбувається ізотропно у власній системі $\pi^0$-мезонів, що рухаються зі швидкістю $v$ відносно лабораторної системи. Знайти відношення числа $\gamma$-квантів, що реєструються у лабораторній системі в передній напівсфері, до числа $\gamma$-квантів у задній напівсфері.

\emph{Застереження:} на відміну від вправи про спостереження зірок, тут фотони рухаються в напрямку від точки розпаду, тому закон перетворення кутів має вигляд $\cos\theta = \frac{\cos\theta' + v/c}{1 + (v/c)\cos\theta'}$ (де штрих позначає власну систему $\pi^0$-мезона).

\emph{Відповідь:} $\frac{1+v/c}{1-v/c}$.
\end{problem}
%=========================================================